% Part: normal-modal-logic
% Chapter: tableaux
% Section: rules-for-K

\documentclass[../../../include/open-logic-section]{subfiles}

\begin{document}

\olfileid{nml}{tab}{rul}

\olsection{قواعد برای~\Log{K}}

قواعد رابط‌های گزاره‌ای معمول همان قواعد !!{tableau}های نشاندار
گزاره‌ای معمول‌اند، تنها با افزودن پیشوندها. در هر مورد، قاعده‌ای که بر
!!{formula} نشاندار~$\sFmla{S}{!A}[\sigma]$ اعمال می‌شود،
!!{formula}های تازه‌ای تولید می‌کند که آن‌ها نیز با~$\sigma$ پیشوند
دارند. قاعده فقط مجاز است این فرمول‌ها را در همان پیشوند تولید کند؛
این ادعا به پیشوند یا جهان دیگری تسری ندارد. این امر باید از نظر شهودی
روشن باشد: برای نمونه، اگر~$!A \land !B$ در (جهان موسوم به)
$\sigma$ صادق باشد، آنگاه~$!A$ و~$!B$ در~$\sigma$ صادق‌اند. قواعد
گزاره‌ای را در~\olref{tab:prop-rules} گرد می‌آوریم.

\begin{table}
  \[\def\arraystretch{3}\begin{array}{|c|c|}
    \hline
    \AxiomC{\sFmla{\True}{\lnot !A}[\sigma]}
    \RightLabel{\TRule{\True}{\lnot}}
    \UnaryInfC{\sFmla{\False}{!A}[\sigma]}
    \DisplayProof
    &
    \AxiomC{\sFmla{\False}{\lnot !A}[\sigma]}
    \RightLabel{\TRule{\False}{\lnot}}
    \UnaryInfC{\sFmla{\True}{!A}[\sigma]}
    \DisplayProof
    \\[1ex]
    \hline
    \AxiomC{\sFmla{\True}{!A \land !B}[\sigma]}
    \RightLabel{\TRule{\True}{\land}}
    \UnaryInfC{\sFmla{\True}{!A}[\sigma]}
    \noLine
    \UnaryInfC{\sFmla{\True}{!B}[\sigma]}
    \DisplayProof
    &
    \AxiomC{\sFmla{\False}{!A \land !B}[\sigma]}
    \RightLabel{\TRule{\False}{\land}}
    \UnaryInfC{$\sFmla{\False}{!A}[\sigma] \quad \mid \quad
      \sFmla{\False}{!B}[\sigma]$}
    \DisplayProof
    \\[2ex]
    \hline
    \AxiomC{\sFmla{\True}{!A \lor !B}[\sigma]}
    \RightLabel{\TRule{\True}{\lor}}
    \UnaryInfC{$\sFmla{\True}{!A}[\sigma] \quad \mid \quad
      \sFmla{\True}{!B}[\sigma]$}
    \DisplayProof
    &
    \AxiomC{\sFmla{\False}{!A \lor !B}[\sigma]}
    \RightLabel{\TRule{\False}{\lor}}
    \UnaryInfC{\sFmla{\False}{!A}[\sigma]}
    \noLine
    \UnaryInfC{\sFmla{\False}{!B}[\sigma]}
    \DisplayProof
    \\[2ex]
    \hline
    \AxiomC{\sFmla{\True}{!A \lif !B}[\sigma]}
    \RightLabel{\TRule{\True}{\lif}}
    \UnaryInfC{$\sFmla{\False}{!A}[\sigma] \quad \mid
      \quad \sFmla{\True}{!B}[\sigma]$}
    \DisplayProof
    &
    \AxiomC{\sFmla{\False}{!A \lif !B}[\sigma]}
    \RightLabel{\TRule{\False}{\lif}}
    \UnaryInfC{\sFmla{\True}{!A}[\sigma]}
    \noLine
    \UnaryInfC{\sFmla{\False}{!B}[\sigma]}
    \DisplayProof
    \\[2ex]
    \hline
  \end{array}\]
  \caption{قواعد !!{tableau} پیشونددار برای رابط‌های گزاره‌ای}
  \ollabel{tab:prop-rules}
\end{table}

شرط بسته‌شدن همانند !!{tableau}های معمول است، هرچند لازم است نه‌فقط
!!{formula}ها بلکه پیشوندها نیز یکسان باشند. پس شاخه‌ای بسته است اگر
هر دو عبارت زیر را شامل شود:
\[
\sFmla{\True}{!A}[\sigma] \quad\text{و}\quad \sFmla{\False}{!A}[\sigma]
\]
برای یک پیشوند~$\sigma$ و !!{formula}~$!A$.

قواعد برپایی فرض‌ها نیز مانند !!{tableau}های معمول است، جز آنکه برای
فرض‌ها همواره از پیشوند~$1$ استفاده می‌کنیم. (مهم نیست کدام پیشوند را
به کار بریم، مشروط بر اینکه برای همهٔ فرض‌ها یکسان باشد.) پس، برای
نمونه، می‌گوییم:
\[
!B_1, \dots, !B_n \Proves !A
\]
اگر و تنها اگر تابلویی بسته برای فرض‌های زیر وجود داشته باشد:
\[
\sFmla{\True}{!B_1}[1], \dots, \sFmla{\True}{!B_n}[1],
\sFmla{\False}{!A}[1].
\]

برای عملگر\iftag{prvBox}{\iftag{prvDiamond}{های موجهاتی~$\Box$
    و~$\Diamond$}{ موجهاتی~$\Box$}}{}\iftag{prvDiamond}{ موجهاتی
    $\Diamond$}{}، پیشوند نتیجهٔ قاعده‌ای که بر !!a{formula} دارای
پیشوند~$\sigma$ اعمال می‌شود، $\sigma.n$ است. بااین‌حال، اینکه کدام
$n$ مجاز است به~$\True$ یا~$\False$ بودن نشانه بستگی دارد.

\iftag{prvBox}{قاعدهٔ~$\TRule{\True}{\Box}$ شاخه‌ای شامل
  $\sFmla{\True}{\Box !A}[\sigma]$ را با
  $\sFmla{\True}{!A}[\sigma.n]$ گسترش می‌دهد.\iftag{prvDiamond}{
    به‌همین‌سان، ق}{}}{ق}\iftag{prvDiamond}{اعدهٔ
  $\TRule{\False}{\Diamond}$ شاخه‌ای شامل
  $\sFmla{\False}{\Diamond !A}[\sigma]$ را با
  $\sFmla{\False}{!A}[\sigma.n]$ گسترش می‌دهد.}{}
\iftag{notprvBox,notprvDiamond}{این قاعده}{این قواعد} را تنها برای
پیشوند~$\sigma.n$ای می‌توان اعمال کرد که \emph{ازپیش} روی شاخهٔ محل
اعمال ظاهر شده باشد. چنین پیشوندی را (روی شاخه) «استفاده‌شده» می‌نامیم.

\iftag{prvBox}{قاعدهٔ~$\TRule{\False}{\Box}$ شاخه‌ای شامل
  $\sFmla{\False}{\Box !A}[\sigma]$ را با
  $\sFmla{\False}{!A}[\sigma.n]$ گسترش می‌دهد.\iftag{prvDiamond}{
    به‌همین‌سان، ق}{}}{ق}\iftag{prvDiamond}{اعدهٔ
  $\TRule{\True}{\Diamond}$ شاخه‌ای شامل
  $\sFmla{\True}{\Diamond !A}[\sigma]$ را با
  $\sFmla{\True}{!A}[\sigma.n]$ گسترش می‌دهد.}{}
\iftag{notprvBox,notprvDiamond}{اما این قاعده}{اما این قواعد} را تنها
برای پیشوند~$\sigma.n$ای می‌توان اعمال کرد که \emph{ازپیش} روی شاخهٔ
محل اعمال ظاهر نشده باشد. چنین پیشوندهایی را برای شاخه «تازه» می‌نامیم.

قواعد در~\olref{tab:rules-K} آمده‌اند.

\begin{table}
  \begin{center}
    \def\arraystretch{3}\def\fCenter{}
    \begin{tabular}{|c|c|}
    \hline
    \iftag{prvBox}{
      \AxiomC{\sFmla{\True}{\Box !A}[\sigma]}
      \RightLabel{\TRule{\True}{\Box}}
      \UnaryInfC{\sFmla{\True}{!A}[\sigma.n]}
      \DisplayProof
      &
      \AxiomC{\sFmla{\False}{\Box !A}[\sigma]}
      \RightLabel{\TRule{\False}{\Box}}
      \UnaryInfC{\sFmla{\False}{!A}[\sigma.n]}
      \DisplayProof\\
      $\sigma.n$ استفاده‌شده است & $\sigma.n$ تازه است
      \\[1ex]
      \hline}{}
    \iftag{prvDiamond}{
      \AxiomC{\sFmla{\True}{\Diamond !A}[\sigma]}
      \RightLabel{\TRule{\True}{\Diamond}}
      \UnaryInfC{\sFmla{\True}{!A}[\sigma.n]}
      \DisplayProof
      &
      \AxiomC{\sFmla{\False}{\Diamond !A}[\sigma]}
      \RightLabel{\TRule{\False}{\Diamond}}
      \UnaryInfC{\sFmla{\False}{!A}[\sigma.n]}
      \DisplayProof\\
      $\sigma.n$ تازه است & $\sigma.n$ استفاده‌شده است
      \\[1ex]
      \hline}{}
    \end{tabular}
  \end{center}
  \caption{قواعد موجهاتی برای~\Log{K}.}
  \ollabel{tab:rules-K}
\end{table}

الزامِ استفاده‌شده‌بودن پیشوند برای
\iftag{prvBox}{$\TRule{\True}{\Box}$}{$\TRule{\False}{\Diamond}$}
ضروری است؛ وگرنه مورد زیر را یک !!{tableau} بسته به شمار می‌آوردیم:
\iftag{notprvBox,notprvDiamond}{%
  \iftag{prvBox}{%
  \begin{oltableau}
    [\pFmla{\True}{\Box \formula{A}}{1}, just = \TAss
      [\pFmla{\False}{\lnot\Box\lnot \formula{A}}{1}, just = \TAss
        [\pFmla{\True}{\formula{A}}{1.1}, just = {\TRule{\True}{\Box}[1]}
          [\pFmla{\True}{\Box\lnot\formula{A}}{1},
            just ={\TRule{\False}{\lnot}[2]}
            [\pFmla{\True}{\lnot\formula{A}}{1.1},
              just = {\TRule{\True}{\Box}[4]}
              [\pFmla{\False}{\formula{A}}{1.1},
                  just ={\TRule{\True}{\lnot}[5]}, close]
            ]
          ]
        ]
      ]
    ]
  \end{oltableau}
  }{
  \begin{oltableau}
    [\pFmla{\True}{\lnot\Diamond\lnot \formula{A}}{1}, just = \TAss
      [\pFmla{\False}{\Diamond\formula{A}}{1}, just = \TAss
        [\pFmla{\False}{\formula{A}}{1.1}, just = {\TRule{\False}{\Diamond}[2]}
          [\pFmla{\False}{\Diamond\lnot\formula{A}}{1},
            just ={\TRule{\True}{\lnot}[1]}
            [\pFmla{\False}{\lnot\formula{A}}{1.1},
              just = {\TRule{\False}{\Diamond}[4]}
              [\pFmla{\True}{\formula{A}}{1.1},
                  just ={\TRule{\True}{\lnot}[5]}, close]
            ]
          ]
        ]
      ]
    ]
  \end{oltableau}
}}{
  \begin{oltableau}
    [\pFmla{\True}{\Box \formula{A}}{1}, just = \TAss
      [\pFmla{\False}{\Diamond \formula{A}}{1}, just = \TAss
        [\pFmla{\True}{\formula{A}}{1.1}, just = {\TRule{\True}{\Box}[1]}
          [\pFmla{\False}{\formula{A}}{1.1},
            just = {\TRule{\False}{\Diamond}[2]}, close]
        ]
      ]
    ]
\end{oltableau}}

اما~$\Box \formula{A} \Entails/ \Diamond \formula{A}$؛ پس دستگاه
اثبات ما نادرست می‌بود. همچنین~$\Diamond \formula{A} \Entails/
\Box \formula{A}$؛ اما بدون الزامِ تازه‌بودن پیشوند برای
\iftag{prvBox}{$\TRule{\False}{\Box}$}{$\TRule{\True}{\Diamond}$}،
مورد زیر تابلویی بسته می‌بود: \iftag{notprvBox,notprvDiamond}{%
  \iftag{prvBox}{%
  \begin{oltableau}
    [\pFmla{\True}{\lnot\Box\lnot \formula{A}}{1}, just = \TAss
      [\pFmla{\False}{\Box \formula{A}}{1}, just = \TAss
        [\pFmla{\False}{\formula{A}}{1.1}, just = {\TRule{\True}{\Box}[2]}
          [\pFmla{\False}{\Box\lnot\formula{A}}{1},
            just ={\TRule{\True}{\lnot}[1]}
            [\pFmla{\False}{\lnot\formula{A}}{1.1},
              just = {\TRule{\False}{\Box}[4]}
              [\pFmla{\True}{\formula{A}}{1.1},
                  just ={\TRule{\False}{\lnot}[5]}, close]
            ]
          ]
        ]
      ]
    ]
  \end{oltableau}
  }{
  \begin{oltableau}
    [\pFmla{\True}{\Diamond \formula{A}}{1}, just = \TAss
      [\pFmla{\False}{\lnot\Diamond\lnot\formula{A}}{1}, just = \TAss
        [\pFmla{\True}{\formula{A}}{1.1}, just = {\TRule{\True}{\Diamond}[1]}
          [\pFmla{\True}{\Diamond\lnot\formula{A}}{1},
            just ={\TRule{\False}{\lnot}[2]}
            [\pFmla{\True}{\lnot\formula{A}}{1.1},
              just = {\TRule{\True}{\Diamond}[4]}
              [\pFmla{\False}{\formula{A}}{1.1},
                  just ={\TRule{\True}{\lnot}[5]}, close]
            ]
          ]
        ]
      ]
    ]
  \end{oltableau}
}}{
  \begin{oltableau}
    [\pFmla{\True}{\Diamond \formula{A}}{1}, just = \TAss,name=one
      [\pFmla{\False}{\Box \formula{A}}{1}, just = \TAss, name=two
        [\pFmla{\True}{\formula{A}}{1.1}, just = {\TRule{\True}{\Diamond}[1]}
          [\pFmla{\False}{\formula{A}}{1.1},
            just = {\TRule{\False}{\Box}[2]}, close]
        ]
      ]
    ]
  \end{oltableau}}

\end{document}
